% GENERATED FILE. Edit canonical lessons, then run npm run generate:books.
% canonical-source-hash: fnv1a64:f07e87075f1ac0e8
% canonical-lessons: MR-C370-vyakhyan, MR-C370-nibandh, MR-C370-bhasan, MR-C370-parisamvad, MR-C370-vakta

\chapter{Talks and Essays}
\label{ch:mr-talks-and-essays}
\hlchaptermodality{370}

\begin{chapteropening}
\textbf{By the end of this chapter:} \emph{I can talk about lectures, speeches, essays and speakers.}

Complete the last lesson of chapter 370: I can talk about lectures, speeches, essays and speakers.
\end{chapteropening}

\section[\texorpdfstring{\mr{व्याख्यान} (vyākhyān)}{व्याख्यान (vyākhyān)}]{\mr{व्याख्यान} (vyākhyān) — a lecture}

\label{lesson:MR-C370-vyakhyan}



\begin{quote}
Before the new one: say the Marathi for a fee, then the Marathi for a hostel.
\end{quote}

\subsection*{You'll want to know: \mr{व्याख्यान}}
\textbf{\mr{व्याख्यान}} — \emph{vyākhyān} — \textquotedblleft{}a lecture\textquotedblright{}.

\begin{quote}
\mr{प्राध्यापकांचं} \mr{व्याख्यान} \mr{दोन} \mr{तास} \mr{चाललं}.
\end{quote}

\begin{quote}
\emph{prādhyāpkāñcaṁ vyākhyān don tās cāllaṁ.}
\end{quote}

\begin{quote}
The professor's lecture went on for two hours.
\end{quote}

Neuter.

\mr{व्याख्यान} \mr{देणे}, to give a lecture.

\begin{cousinweb}
\textbf{\mr{व्याख्यान}} (\emph{vyākhyān}) is a tatsama, taken from Sanskrit \emph{vyākhyāna}, \textquotedblleft{}comment, narration\textquotedblright{} (CDIAL 12188); the same entry lists an inherited Marathi \emph{vākhāṇ}, \textquotedblleft{}praise\textquotedblright{}.
\end{cousinweb}

\begin{grammarlens}[title={What you've built}]
One more word for giving reasons and opinions.
\end{grammarlens}

\subsection*{Your turn}
\begin{itemize}
\raggedright
  \item \emph{Say it:} \emph{vyākhyān}
  \item \emph{Say it:} \emph{vyākhyān}, once more
  \item \emph{Say it:} \emph{prādhyāpkāñcaṁ vyākhyān don tās cāllaṁ.} — The professor's lecture went on for two hours.
  \item \emph{Recall:} say the Marathi for a fee, then the Marathi for a hostel, then say \emph{vyākhyān} again
\end{itemize}

\begin{tcolorbox}[breakable,colback=teal!4,colframe=teal!35!black,title={Before you move on}]
What does \mr{व्याख्यान} mean? (\textquotedblleft{}A lecture\textquotedblright{}.) Say it once more.
\end{tcolorbox}
\section[\texorpdfstring{\mr{निबंध} (nibandh)}{निबंध (nibandh)}]{\mr{निबंध} (nibandh) — an essay}

\label{lesson:MR-C370-nibandh}



\begin{quote}
Before the new one: say the Marathi for a hostel, then the Marathi for a lecture.
\end{quote}

\subsection*{You'll want to know: \mr{निबंध}}
\textbf{\mr{निबंध}} — \emph{nibandh} — \textquotedblleft{}an essay\textquotedblright{}.

\begin{quote}
\mr{मी} \mr{पावसाळ्यावर} \mr{निबंध} \mr{लिहिला}.
\end{quote}

\begin{quote}
\emph{mī pāvsāḷyāvar nibandh lihilā.}
\end{quote}

\begin{quote}
I wrote an essay on the rainy season.
\end{quote}

Masculine.

X \mr{वर} \mr{निबंध}, an essay on X.

\begin{cousinweb}
\textbf{\mr{निबंध}} (\emph{nibandh}) is a tatsama, taken from Sanskrit \emph{nibandha}, \textquotedblleft{}fastening; a literary composition\textquotedblright{} (CDIAL 7222).
\end{cousinweb}

\begin{grammarlens}[title={What you've built}]
One more word for giving reasons and opinions.
\end{grammarlens}

\subsection*{Your turn}
\begin{itemize}
\raggedright
  \item \emph{Say it:} \emph{nibandh}
  \item \emph{Say it:} \emph{nibandh}, once more
  \item \emph{Say it:} \emph{mī pāvsāḷyāvar nibandh lihilā.} — I wrote an essay on the rainy season.
  \item \emph{Recall:} say the Marathi for a hostel, then the Marathi for a lecture, then say \emph{nibandh} again
\end{itemize}

\begin{tcolorbox}[breakable,colback=teal!4,colframe=teal!35!black,title={Before you move on}]
What does \mr{निबंध} mean? (\textquotedblleft{}An essay\textquotedblright{}.) Say it once more.
\end{tcolorbox}
\section[\texorpdfstring{\mr{भाषण} (bhāṣaṇ)}{भाषण (bhāṣaṇ)}]{\mr{भाषण} (bhāṣaṇ) — a speech}

\label{lesson:MR-C370-bhasan}



\begin{quote}
Before the new one: say the Marathi for a lecture, then the Marathi for an essay.
\end{quote}

\subsection*{You'll want to know: \mr{भाषण}}
\textbf{\mr{भाषण}} — \emph{bhāṣaṇ} — \textquotedblleft{}a speech\textquotedblright{}.

\begin{quote}
\mr{मुख्य} \mr{पाहुण्यांनी} \mr{छोटं} \mr{भाषण} \mr{केलं}.
\end{quote}

\begin{quote}
\emph{mukhya pāhuṇyānnī choṭaṁ bhāṣaṇ kelaṁ.}
\end{quote}

\begin{quote}
The chief guest gave a short speech.
\end{quote}

Neuter.

\mr{भाषण} \mr{करणे} or \mr{देणे}, to make a speech.

\begin{cousinweb}
\textbf{\mr{भाषण}} (\emph{bhāṣaṇ}) is a tatsama, taken from Sanskrit \emph{bhāṣaṇa}, \textquotedblleft{}speaking, talk\textquotedblright{} (Monier-Williams).
\end{cousinweb}

\begin{grammarlens}[title={What you've built}]
One more word for giving reasons and opinions.
\end{grammarlens}

\subsection*{Your turn}
\begin{itemize}
\raggedright
  \item \emph{Say it:} \emph{bhāṣaṇ}
  \item \emph{Say it:} \emph{bhāṣaṇ}, once more
  \item \emph{Say it:} \emph{mukhya pāhuṇyānnī choṭaṁ bhāṣaṇ kelaṁ.} — The chief guest gave a short speech.
  \item \emph{Recall:} say the Marathi for a lecture, then the Marathi for an essay, then say \emph{bhāṣaṇ} again
\end{itemize}

\begin{tcolorbox}[breakable,colback=teal!4,colframe=teal!35!black,title={Before you move on}]
What does \mr{भाषण} mean? (\textquotedblleft{}A speech\textquotedblright{}.) Say it once more.
\end{tcolorbox}
\section[\texorpdfstring{\mr{परिसंवाद} (parisaṁvād)}{परिसंवाद (parisaṁvād)}]{\mr{परिसंवाद} (parisaṁvād) — a symposium; a panel discussion}

\label{lesson:MR-C370-parisamvad}



\begin{quote}
Before the new one: say the Marathi for an essay, then the Marathi for a speech.
\end{quote}

\subsection*{You'll want to know: \mr{परिसंवाद}}
\textbf{\mr{परिसंवाद}} — \emph{parisaṁvād} — \textquotedblleft{}a symposium; a panel discussion\textquotedblright{}.

\begin{quote}
\mr{महाविद्यालयात} \mr{पर्यावरणावर} \mr{परिसंवाद} \mr{झाला}.
\end{quote}

\begin{quote}
\emph{mahāvidyālyāt paryāvarṇāvar parisaṁvād jhālā.}
\end{quote}

\begin{quote}
A symposium on the environment was held at the college.
\end{quote}

Masculine.

\mr{परि} + \mr{संवाद}, around + conversation.

\begin{cousinweb}
\textbf{\mr{परिसंवाद}} (\emph{parisaṁvād}) is a learned word built from Sanskrit \emph{pari-}, \textquotedblleft{}around\textquotedblright{}, and \emph{saṁvāda}, \textquotedblleft{}conversation\textquotedblright{} (Monier-Williams).
\end{cousinweb}

\begin{grammarlens}[title={What you've built}]
One more word for giving reasons and opinions.
\end{grammarlens}

\subsection*{Your turn}
\begin{itemize}
\raggedright
  \item \emph{Say it:} \emph{parisaṁvād}
  \item \emph{Say it:} \emph{parisaṁvād}, once more
  \item \emph{Say it:} \emph{mahāvidyālyāt paryāvarṇāvar parisaṁvād jhālā.} — A symposium on the environment was held at the college.
  \item \emph{Recall:} say the Marathi for an essay, then the Marathi for a speech, then say \emph{parisaṁvād} again
\end{itemize}

\begin{tcolorbox}[breakable,colback=teal!4,colframe=teal!35!black,title={Before you move on}]
What does \mr{परिसंवाद} mean? (\textquotedblleft{}A symposium; a panel discussion\textquotedblright{}.) Say it once more.
\end{tcolorbox}
\section[\texorpdfstring{\mr{वक्ता} (vaktā)}{वक्ता (vaktā)}]{\mr{वक्ता} (vaktā) — a speaker}

\label{lesson:MR-C370-vakta}



\begin{quote}
Before the new one: say the Marathi for a speech, then the Marathi for a symposium.
\end{quote}

\subsection*{You'll want to know: \mr{वक्ता}}
\textbf{\mr{वक्ता}} — \emph{vaktā} — \textquotedblleft{}a speaker\textquotedblright{}.

\begin{quote}
\mr{आजचे} \mr{वक्ते} \mr{पुण्याहून} \mr{आले} \mr{आहेत}.
\end{quote}

\begin{quote}
\emph{ājce vakte puṇyāhūn āle āhet.}
\end{quote}

\begin{quote}
Today's speaker has come from Pune.
\end{quote}

Masculine. Plural: \mr{वक्ते}, \emph{vakte}.

Plural \mr{वक्ते}; respect takes the plural; a listener is \mr{श्रोता}.

\begin{cousinweb}
\textbf{\mr{वक्ता}} (\emph{vaktā}) is a tatsama, taken from Sanskrit \emph{vaktṛ}, \textquotedblleft{}a speaker\textquotedblright{} (Monier-Williams). The Sanskrit nominative is \emph{vaktā}.
\end{cousinweb}

\begin{grammarlens}[title={What you've built}]
One more word for giving reasons and opinions.
\end{grammarlens}

\subsection*{Your turn}
\begin{itemize}
\raggedright
  \item \emph{Say it:} \emph{vaktā}
  \item \emph{Say it:} \emph{vaktā}, once more
  \item \emph{Say it:} \emph{ājce vakte puṇyāhūn āle āhet.} — Today's speaker has come from Pune.
  \item \emph{Recall:} say the Marathi for a speech, then the Marathi for a symposium, then say \emph{vaktā} again
\end{itemize}

\begin{tcolorbox}[breakable,colback=teal!4,colframe=teal!35!black,title={Before you move on}]
What does \mr{वक्ता} mean? (\textquotedblleft{}A speaker\textquotedblright{}.) Say it once more.
\end{tcolorbox}
